\section{Métodos}

\subsection{Exercício 1: FFT}

Foi criada uma matriz aleatória \texttt{float32} de dimensão
$3000\times3000$. A FFT foi aplicada sobre a primeira dimensão, reproduzindo o
comportamento do exemplo do enunciado. Foram medidos:

\begin{enumerate}
    \item FFT em CPU com NumPy;
    \item somente a FFT em GPU;
    \item transferência host--device seguida da FFT;
    \item custo médio por FFT ao reutilizar dados residentes na GPU.
\end{enumerate}

Eventos CUDA e sincronização explícita foram usados para impedir que a
execução assíncrona produzisse tempos incorretos. Os resultados CPU e GPU
foram comparados com tolerância numérica.

\subsection{Exercício 2a: CPU}

O plano complexo $[-2,1]\times[-1{,}5,1{,}5]$ foi mapeado para uma matriz
$1000\times1000$. Para cada pixel, o algoritmo:

\begin{enumerate}
    \item calcula o número complexo $c$ correspondente;
    \item inicializa $Z_0=c$;
    \item aplica a recorrência até 500 vezes;
    \item interrompe quando $|Z|>2$;
    \item grava na matriz a última iteração válida.
\end{enumerate}

A versão serial foi compilada com \texttt{numba.njit}.

\subsection{Exercício 2b: GPU vetorizada}

A alternativa ao \texttt{gpuArray} do MATLAB foi implementada com matrizes
CuPy. A cada iteração, operações vetorizadas atualizam todos os pontos ativos
da matriz simultaneamente.

\subsection{Exercício 2c: kernel CUDA}

A alternativa ao \texttt{arrayfun} foi um kernel Numba CUDA com um thread por
pixel \autocite{numba}. Foram usados blocos de $16\times16$ threads. Todas as
versões utilizaram precisão \texttt{float32}, garantindo uma comparação
direta das matrizes resultantes.
